\chapter{Przestrzenie klasyfikujące i teoria przeszkód}
\label{ch:klasyfikacja-v6}

\section{Dlaczego potrzebujemy przestrzeni klasyfikujących?}
W rozdziałach o homologii i kohomologii przypisaliśmy przestrzeni $X$
grupy $H^n(X;A)$. Czy grupy te można opisać za pomocą
\emph{odwzorowań} z $X$ do jednej przestrzeni, zależnej tylko od
$A$ i $n$? Podobnie: czy wszystkie wiązki danego typu nad $X$
można uzyskać jako cofnięcia jednej wiązki uniwersalnej?
Na oba pytania odpowiedź jest twierdząca po sprecyzowaniu
hipotez. Wprowadzimy przestrzenie Eilenberga--Mac Lane'a,
przestrzenie klasyfikujące grup i teorię przeszkód, która
wyjaśnia, dlaczego obiekty globalne mogą nie istnieć
pomimo lokalnego istnienia.

Zakładamy w tym rozdziale, że przestrzenie służące za
podstawy klasyfikacji są CW-kompleksami. Gdy mówimy
o ogólnych wiązkach głównych nad przestrzenią $X$,
zakładamy, że $X$ jest parazwarta, a wiązka jest
\emph{numerowalna} (istnieje podział jedności podporządkowany
pokryciu trywializującemu). Dla CW-kompleksów
warunek ten nie stanowi dodatkowego ograniczenia
w standardowej teorii wiązek lokalnie trywialnych.

\section{Przestrzenie Eilenberga--Mac Lane'a}
\begin{definicja}
Niech $n\ge1$. Dla $n=1$ niech $G$ będzie dowolną grupą;
dla $n\ge2$ niech $G$ będzie grupą abelową.
\emph{Przestrzenią Eilenberga--Mac Lane'a typu $K(G,n)$}
nazywamy łukowo spójny CW-kompleks $K$ z wyróżnionym
punktem, dla którego
\[
\pi_j(K)=\begin{cases}G,&j=n,\\0,&j\ge1,\ j\ne n.
\end{cases}
\]
Przy identyfikacji $\pi_n(K)\cong G$ ustalamy
izomorfizm, a nie tylko istnienie izomorfizmu.
\end{definicja}
\begin{przyklad}
Mamy $S^1\simeq K(\Z,1)$, a nieskończeniewymiarowa
przestrzeń rzutowa rzeczywista $\R P^\infty$ jest
$K(\Z/2,1)$: jej dwukrotne nakrycie $S^\infty$
jest ściągalne. Pokażemy także, że
$\C P^\infty\simeq K(\Z,2)$.
Sfera $S^n$ dla $n\ge2$ \emph{nie} jest przestrzenią
$K(\Z,n)$, gdyż posiada niezerowe grupy homotopii
w wyższych wymiarach.
\end{przyklad}

\begin{twierdzenie}[Istnienie przestrzeni $K(G,n)$]
\label{thm:KGn-existence-v6}
Dla każdej grupy $G$ istnieje CW-kompleks $K(G,1)$.
Dla każdej grupy abelowej $A$ i $n\ge2$ istnieje
CW-kompleks $K(A,n)$. Każdy taki obiekt jest jednoznaczny
z dokładnością do równoważności homotopijnej zachowującej
ustaloną identyfikację odpowiedniej grupy homotopii.
\end{twierdzenie}
\begin{proof}
Najpierw skonstruujmy $K(G,1)$. Wybierzmy prezentację
grupy $G$ przez generatory i relacje, niekoniecznie skończoną.
Zbudujmy spójny dwuwymiarowy CW-kompleks $X_2$:
jeden wierzchołek, jedna $1$-komórka dla każdego generatora
i jedna $2$-komórka przyklejona wzdłuż słowa odpowiadającego
każdej relacji. Twierdzenie van Kampena daje $\pi_1(X_2)\cong G$.
Następnie dla wszystkich generatorów grupy $\pi_2(X_2)$
doklejamy $3$-komórki wzdłuż reprezentujących je map
$S^2\to X_2$. W ten sposób zabijamy $\pi_2$, nie zmieniając
$\pi_1$. Powtarzamy indukcyjnie: na etapie $k$ doklejamy
$(k+1)$-komórki zabijające $\pi_k$, bez zmiany $\pi_j$
dla $j<k$. W słabej topologii granicznego CW-kompleksu
każda mapa sfery ma obraz homotopijnie reprezentowalny
w skończonym etapie odpowiedniej filtracji komórkowej;
stąd otrzymujemy $\pi_k=0$ dla $k\ge2$.

Dla grupy abelowej $A$ i $n\ge2$ wybieramy prezentację
$0\to R\to F\to A\to0$, w której $R$ i $F$ są wolne.
Bierzemy bukiet sfer $S^n$ indeksowany bazą $F$;
reprezentanci relacji $R\subset F\cong\pi_n(\bigvee S^n)$
wyznaczają mapy dołączające $(n+1)$-komórki.
Powstały kompleks jest $(n-1)$-spójny i ma
$H_n\cong A$, zatem z twierdzenia Hurewicza
$\pi_n\cong A$. Dołączając kolejno komórki wymiarów
$n+2,n+3,\ldots$ zabijamy pozostałe grupy homotopii,
nie zmieniając $\pi_n$. Otrzymujemy $K(A,n)$.

Aby wyjaśnić jednoznaczność, zauważmy, że dla dwóch takich
kompleksów $K,L$ homomorfizm $\pi_n(K)\to\pi_n(L)$
można zrealizować odwzorowaniem komórkowym: na
$n$-szkielecie podajemy reprezentanty grupy $\pi_n(L)$,
a przeszkody rozszerzania na wyższe komórki znikają,
ponieważ $\pi_j(L)=0$ dla $j>n$. Dla $n=1$ realizację
homomorfizmu $\pi_1$ uzyskujemy z prezentacji przez
wypełnienie relacji, a dalsze komórki dołączamy, gdyż
$\pi_j(L)=0$ dla $j>1$. W ten sposób konstruujemy mapę
indukującą zadany izomorfizm $\pi_n$ (odpowiednio $\pi_1$).
Jest to słaba równoważność homotopijna, a więc na mocy
twierdzenia Whiteheada równoważność homotopijna.
\end{proof}

\section{Kohomologia jako zbiór klas homotopii}
Niech $[X,Y]$ oznacza zbiór klas homotopii ciągłych
odwzorowań $X\to Y$. Dla $n\ge1$ i grupy abelowej $A$
zbiór $[X,K(A,n)]$ ma naturalną strukturę grupy abelowej,
choć nie jest ona oczywista z samej definicji homotopii.

\begin{twierdzenie}[Reprezentacja kohomologii singularnej]
\label{thm:rep-coho-v6}
Dla dowolnego CW-kompleksu $X$, grupy abelowej $A$ i
$n\ge1$ istnieje naturalna bijekcja
\[
[X,K(A,n)]\cong H^n(X;A).
\]
Jest ona izomorfizmem grup przy naturalnej strukturze
na zbiorze klas homotopii po lewej stronie.
Dla $n\ge2$ tę samą bijekcję można sformułować
w wersji z wyróżnionym punktem dla spójnych CW-kompleksów.
\end{twierdzenie}
\begin{proof}
Pokażemy szczegółowo mechanizm komórkowy dla $n\ge2$;
przypadek $n=1$ otrzymuje się analogicznie z $K(A,1)$
lub przez klasyfikację homomorfizmów do grupy abelowej $A$.

\emph{Krok 1: kołańcuch z odwzorowania.}
Ponieważ $K(A,n)$ jest $(n-1)$-spójna, za pomocą
przybliżenia komórkowego i kolejnych homotopii na komórkach
możemy doprowadzić mapę $f:X\to K(A,n)$ do postaci
stałej na $X^{n-1}$. Na każdej zorientowanej $n$-komórce
$e^n_\alpha$ mapa ilorazowa
$D^n/\partial D^n\cong S^n\to K(A,n)$
wyznacza element $a_\alpha\in\pi_n(K(A,n))=A$.
Przypisanie $e^n_\alpha\mapsto a_\alpha$ definiuje
kołańcuch komórkowy $c_f\in C^n_{\mathrm{CW}}(X;A)$.

\emph{Krok 2: warunek kocyklu.}
Niech $e^{n+1}$ będzie komórką dołączoną mapą
$\varphi:S^n\to X^n$. Obstrukcja rozszerzenia mapy
z brzegu na jej dysk leży w $\pi_n(K(A,n))=A$.
Po zastosowaniu komórkowego przybliżenia $\varphi$
obstrukcja jest sumą elementów $a_\alpha$ ważonych
liczbami incydencji $[e^{n+1}:e^n_\alpha]$.
Jest to dokładnie $(\delta c_f)(e^{n+1})$.
Ponieważ $f$ istnieje na całym $X$, otrzymujemy
$\delta c_f=0$.

\emph{Krok 3: realizacja dowolnego kocyklu.}
Odwrotnie, dla $c\in Z^n_{\mathrm{CW}}(X;A)$
ustalamy mapę stałą na $X^{n-1}$ i wybieramy na
każdej $n$-komórce reprezentanta elementu $c(e^n)$.
Warunek $\delta c=0$ jest dokładnie warunkiem
wypełnienia map na granicach wszystkich $(n+1)$-komórek.
Na komórkach wyższych wymiarów przeszkody leżą
w $\pi_j(K(A,n))$ dla $j>n$ i dlatego znikają.
Powstaje mapa $f_c:X\to K(A,n)$.

\emph{Krok 4: homotopie i kobrzegi.}
Jeżeli $c'=c+\delta b$, gdzie
$b\in C^{n-1}_{\mathrm{CW}}(X;A)$, konstruujemy
homotopię na $X\times I$ przez nadanie na pryzmatach
$(n-1)$-komórek elementów $b(e^{n-1})\in A$.
Warunek możliwości rozszerzenia przez pryzmaty
$n$-komórek ma postać
$c'-c-\delta b=0$; powyżej tych wymiarów obstrukcje
znikają. Odwrotnie każda homotopia $f\simeq g$ po
komórkowym uporządkowaniu względem $X\times\partial I$
wyznacza kołańcuch $b$ z $c_g-c_f=\delta b$.
Zatem klasy homotopii odpowiadają dokładnie klasom
$Z^n/B^n=H^n_{\mathrm{CW}}(X;A)$, które są naturalnie
izomorficzne z kohomologią singularną.
Dodawanie klas odpowiada dodawaniu reprezentujących
je kocykli, co nadaje lewej stronie strukturę grupową.
\end{proof}

\begin{przyklad}[Dlaczego samą grupą nie da się zastąpić pierścienia]
Niech $h\in H^2(\C P^2;\Z)$ będzie generatorem.
Z twierdzenia reprezentacyjnego odpowiada mu mapa
$f:\C P^2\to K(\Z,2)$. Iloczyn
$h\smile h\in H^4(\C P^2;\Z)$ jest niezerowy,
natomiast na $S^2\vee S^4$ iloczyn dwóch klas
stopnia $2$ jest zerowy. Grupy addytywne homologii
obu przestrzeni są jednak izomorficzne. Stąd pełny
pierścień kohomologii zawiera informację niewidoczną
w samych grupach.
\end{przyklad}

\section{Przestrzeń klasyfikująca grupy dyskretnej}
\begin{definicja}[Wiązka główna]
Niech $G$ będzie grupą topologiczną. \emph{Główną
wiązką $G$} nad $B$ nazywamy lokalnie trywialne
rzutowanie $P\to B$ z ciągłym, wolnym prawym działaniem
$G$ na $P$, które lokalnie identyfikuje $P|_U$
z $U\times G$ i działanie z prawym mnożeniem na $G$.
Izomorfizmy wiązek mają zachowywać zarówno rzutowanie,
jak i działanie $G$.
\end{definicja}

\begin{twierdzenie}[Konstrukcja $EG\to BG$ dla grupy dyskretnej]
\label{thm:bar-BG-v6}
Dla dowolnej grupy dyskretnej $G$ istnieje
ściągalny CW-kompleks $EG$ z wolnym działaniem $G$,
dla którego iloraz $BG=EG/G$ jest $K(G,1)$.
Rzutowanie $EG\to BG$ jest uniwersalną główną
wiązką $G$ wśród wiązek nad CW-kompleksami.
\end{twierdzenie}
\begin{proof}
Konstrukcję można podać bez odwołania do
nieskończonych przestrzeni Grassmanna. Bierzemy
realizację geometryczną zbioru symplicjalnego
złożonego z ciągów $(g_0,\ldots,g_k)\in G^{k+1}$:
opuszczanie jednego wyrazu definiuje ściany,
a powtarzanie wyrazu degeneracje. Równoczesne
lewe mnożenie wszystkich współrzędnych daje wolne
działanie $G$ na realizacji $EG$. Operator dodający
z przodu element neutralny,
$(g_0,\ldots,g_k)\mapsto(1,g_0,\ldots,g_k)$,
jest dodatkową degeneracją; odpowiadające mu stożkowanie
w realizacji daje kontrakcję $EG$.
Rzutowanie $EG\to BG$ jest nakryciem o grupie
transformacji pokrycia $G$ i ma ściągalną przestrzeń
całkowitą, zatem $\pi_1(BG)\cong G$ i
$\pi_j(BG)=0$ dla $j>1$. Jest więc modelem $K(G,1)$.

Dla spójnego CW-kompleksu $X$ główne wiązki
grupy dyskretnej $G$ odpowiadają homomorfizmom
$\pi_1(X)\to G$ z dokładnością do sprzężenia:
homomorfizm jest monodromią nakrycia, a nakrycie
odtwarzamy jako iloraz $\widetilde X\times G$
przez działanie grupy $\pi_1(X)$ wyznaczone tym
homomorfizmem. Z drugiej strony
$[X,BG]$ jest tym samym zbiorem klas homomorfizmów,
ponieważ $BG=K(G,1)$. Zatem każdą wiązkę otrzymujemy
jako cofnięcie $EG\to BG$; zgodność z homotopiami
map klasyfikujących wynika z niezmienniczości
monodromii względem homotopii.
Dla niespójnego $X$ konstrukcję stosujemy oddzielnie
do każdej składowej.
\end{proof}

\begin{uwaga}[Nienaturalność wyboru monodromii]
Po ustaleniu punktów w bazie i we włóknie otrzymuje się
konkretny homomorfizm $\pi_1(X)\to G$.
Bez wyboru punktu we włóknie zmiana trywializacji
sprzęga ten homomorfizm przez element $G$.
Dlatego w klasyfikacji \emph{niepunktowanych} wiązek
występują klasy sprzężoności.
\end{uwaga}

\section{Grupy topologiczne, wiązki uniwersalne i przestrzenie Grassmanna}
Dla grupy topologicznej $G$ model $BG$ można otrzymać
z nieskończonego złączenia $EG=G*G*G*\cdots$.
Punkt w tym złączeniu ma postać skończonej formalnej
kombinacji $\sum_i t_i g_i$, gdzie $t_i\ge0$,
$\sum_i t_i=1$, a współrzędną $g_i$ pomijamy,
gdy $t_i=0$. Grupa $G$ działa równocześnie prawym
mnożeniem na wszystkich niezerowych współrzędnych.

\begin{lemat}[Ściągalność modelu uniwersalnego]
Nieskończone złączenie $EG=G*G*\cdots$ jest ściągalne,
a diagonalne działanie $G$ jest wolne.
\end{lemat}
\begin{proof}
Każdy punkt ma co najmniej jedną współrzędną
o dodatniej wadze, na której prawostronna translacja
jest wolna; zatem całe działanie jest wolne.
Dla ściągalności wygodny jest równoważny
\emph{symplicjalny model} przestrzeni $EG$ z
$k$-sympleksami $G^{k+1}$, którego operatory ścian
opuszczają współrzędne, a degeneracje je powtarzają.
Operacja
$(g_0,\ldots,g_k)\mapsto(1_G,g_0,\ldots,g_k)$
jest dodatkową degeneracją, a odpowiadający jej
stożek na realizacji daje homotopię identyczności
do mapy stałej. Przejście pomiędzy tym modelem
a nieskończonym złączeniem zachowuje typ homotopijny
(oba są standardowymi modelami uniwersalnej
numerowalnej wiązki); to wystarcza do wniosku.
\end{proof}

\begin{twierdzenie}[Klasyfikacja numerowalnych wiązek głównych]
\label{thm:BG-general-v6}
Niech $G$ będzie grupą topologiczną, $X$ przestrzenią
parazwartą. Klasy izomorfizmu numerowalnych głównych
wiązek $G$ nad $X$ pozostają w naturalnej bijekcji z
$[X,BG]$, gdzie $BG=EG/G$ dla uniwersalnego
numerowalnego modelu $EG\to BG$.
\end{twierdzenie}
\begin{proof}
Pokażmy najpierw, skąd bierze się mapa klasyfikująca.
Niech $P\to X$ ma pokrycie trywializujące $\{U_i\}$
z lokalnymi przekrojami $s_i$, a $\{\lambda_i\}$ będzie
podporządkowanym mu lokalnie skończonym podziałem jedności.
Dla $e\in P_x$ piszemy $e=s_i(x)g_i(e)$ na $U_i$.
Przypisanie
\[
\Phi(e)=\sum_i\lambda_i(x)g_i(e)\in EG
\]
jest ciągłe (lokalnie występuje skończenie wiele składników)
i $G$-ekwawariantne: $\Phi(eh)=\Phi(e)h$.
Przechodząc do ilorazu otrzymujemy mapę
$f:X\to BG$ oraz izomorfizm wiązek
$P\cong f^*(EG)$. To sprawdza się w każdej trywializacji:
współrzędna $g_i(e)$ przekształca się przez prawą
translację, więc na włóknach otrzymujemy bijekcję
$G$-przestrzeni zgodną z topologią.

Dwa podziały jedności i dwa wybory lokalnych przekrojów
prowadzą do homotopijnych map klasyfikujących:
umieszczamy wybrane konstrukcje w rozłącznych
zbiorach współrzędnych nieskończonego złączenia
i interpolujemy współczynniki przez $(1-t)$ oraz $t$.
Analogiczna interpolacja pokazuje, że izomorficzne
wiązki wyznaczają homotopijne mapy. Odwrotnie,
cofnięcia wiązki uniwersalnej wzdłuż map homotopijnych
są izomorficzne, ponieważ numerowalne wiązki nad
$X\times I$ mają izomorficzne restrykcje do końców
cylindra. Zbieżność konstrukcji nie wymaga nieskończonej
sumy w otoczeniu pojedynczego punktu, gdyż użyliśmy
lokalnie skończonego podziału jedności.
\end{proof}

\begin{center}
\begin{tikzpicture}[>=Stealth,node distance=1.7cm]
\node (P) {$P$};
\node (EG) [right=3.4cm of P] {$EG$};
\node (X) [below=1.4cm of P] {$X$};
\node (BG) [below=1.4cm of EG] {$BG$};
\draw[->,thick] (P)--(EG) node[midway,above]{$\Phi$};
\draw[->,thick] (P)--(X) node[midway,left]{$\pi$};
\draw[->,thick] (EG)--(BG) node[midway,right]{$\pi_{\rm univ}$};
\draw[->,thick] (X)--(BG) node[midway,below]{$f$};
\end{tikzpicture}
\end{center}
Powyższy diagram jest przemienny, a górna mapa
$\Phi$ jest $G$-ekwawariantna. Taką własność ma
mapa klasyfikująca skonstruowana w dowodzie.

\begin{wniosek}[Klasyfikacja wiązek wektorowych]
Niech $X$ będzie parazwartą przestrzenią, $r\ge1$.
Izomorfizmy rzeczywistych wiązek wektorowych rzędu $r$
nad $X$ odpowiadają klasom homotopii map
$X\to BO(r)$, a zespolonych rzędu $r$ -- klasom map
$X\to BU(r)$. Standardowe modele mają postać
granic przestrzeni Grassmanna skończeniewymiarowych.
\end{wniosek}
\begin{proof}
Metryka włóknowa na rzeczywistej wiązce redukuje
grupę strukturalną do $O(r)$. Jej wiązka ram
ortonormalnych jest główną wiązką $O(r)$,
a wiązkę pierwotną odzyskujemy przez dołączenie
reprezentacji standardowej na $\R^r$.
Stosujemy Twierdzenie~\ref{thm:BG-general-v6}.
W przypadku zespolonym wybieramy metrykę hermitowską
i zastępujemy $O(r)$ przez $U(r)$.
\end{proof}

\begin{przyklad}[Zespolone wiązki liniowe i $\C P^\infty$]
Grupa $U(1)=S^1$ działa swobodnie na
$S^\infty=\bigcup_{N\ge0}S^{2N+1}$, który jest
ściągalny. Iloraz to $\C P^\infty$, zatem
$BU(1)\simeq\C P^\infty$. Z długiego ciągu homotopii
fibracji $S^1\to S^\infty\to\C P^\infty$ otrzymujemy
$\pi_2(\C P^\infty)\cong\Z$ i
$\pi_j(\C P^\infty)=0$ dla $j\ne2$, $j\ge1$.
Stąd $\C P^\infty=K(\Z,2)$ i dla CW-kompleksu $X$
\[
\{\text{zespolone wiązki liniowe nad }X\}/\cong
\quad\longleftrightarrow\quad H^2(X;\Z).
\]
Odpowiadająca klasie kohomologii liczba lub klasa to
pierwsza klasa Cherna $c_1(L)$. W szczególności
nad $S^2$ wiązki liniowe klasyfikuje liczba całkowita
$\langle c_1(L),[S^2]\rangle$.
\end{przyklad}

\section{Teoria przeszkód: rozszerzanie map i przekrojów}
\begin{definicja}[Przeszkoda rozszerzania]
Niech $(X,A)$ będzie parą CW, a $Y$ przestrzenią
łukowo spójną. Załóżmy, że mapa $f$ została określona
na $A\cup X^n$. Dla $(n+1)$-komórki
$e_\alpha^{n+1}$ z mapą charakterystyczną
$\chi_\alpha:D^{n+1}\to X$ ograniczenie
$f\circ\chi_\alpha|_{S^n}$ reprezentuje klasę
$[f\circ\chi_\alpha|_{S^n}]\in\pi_n(Y)$,
po przetransportowaniu punktu bazowego.
Rodzinę tych klas nazywamy \emph{kołańcuchem przeszkody}
rozszerzenia na komórki stopnia $n+1$.
Dla $n\ge2$ jest to kołańcuch o wartościach
w systemie lokalnym otrzymanym z działania
$\pi_1(Y)$ na $\pi_n(Y)$, cofniętym przez
$f_*:\pi_1(X)\to\pi_1(Y)$ (mapa $f$ jest już
określona na $2$-szkielecie). Transporty wymagają
ustalenia dróg łączących punkty bazowe. Jeżeli
to działanie jest trywialne, współczynniki są stałe.
Dla pary $(X,A)$ w rachunku względnym uwzględniamy
te same transporty w $X$.
\end{definicja}

\begin{center}
\begin{tikzpicture}[>=Stealth]
 \draw[thick,fill=blue!4] (0,0) circle (.88);
 \draw[very thick,red!70!black] (0,0) circle (.88);
 \node at (0,0){$D^{n+1}$};
 \node[red!70!black,below] at (0,-.95){$S^n$};
 \draw[->,thick] (1.2,.15)--(3.45,.15)
   node[midway,above]{$f\circ\chi|_{S^n}$};
 \draw[thick,fill=blue!4] (4.4,0) ellipse (.7 and .85);
 \node at (4.4,0){$Y$};
 \node[align=center] at (2.3,-1.78){\footnotesize Przedłużenie na dysk istnieje dokładnie wtedy,\\
 \footnotesize gdy klasa mapy na brzegu w $\pi_n(Y)$ jest zerowa.};
\end{tikzpicture}
\end{center}

\begin{twierdzenie}[Pierwsza przeszkoda dla sekcji]
\label{thm:obstruction-section-v6}
Niech $p:E\to B$ będzie lokalnie trywialną wiązką nad
spójnym CW-kompleksem, której włókno $F$ jest
$(n-1)$-spójne dla pewnego $n\ge2$.
Przypuśćmy, że skonstruowano przekrój $s$ nad
$n$-szkieletem $B^n$. Dołączenie każdej $(n+1)$-komórki
wyznacza element $\pi_n(F)$, a elementy te tworzą
kocykl komórkowy $o_{n+1}(s)$ o wartościach
w systemie lokalnym $\mathcal\pi_n(F)$.
Przy ustalonym przekroju nad $(n-1)$-szkieletem
jego klasa w $H^{n+1}(B;\mathcal\pi_n(F))$ nie zależy
od jego rozszerzenia na $n$-szkielet. Ponieważ
włókno jest $(n-1)$-spójne, dwa przekroje nad
$(n-1)$-szkieletem można połączyć odpowiednią
homotopią; wynikająca stąd klasa jest więc
niezależna od wyborów (po identyfikacji systemów).
Nazywa się ją \emph{pierwszą przeszkodą} istnienia przekroju. Znikanie tej klasy pozwala
skonstruować przekrój na $B^{n+1}$ (po ewentualnej
zmianie przekroju na $B^n$), lecz nie zapewnia
jeszcze jego istnienia na całym $B$, jeżeli
wyższe grupy $\pi_j(F)$ są niezerowe.
\end{twierdzenie}
\begin{proof}
Cofnięcie wiązki na dysk charakterystyczny każdej
$(n+1)$-komórki jest trywialne. Przekrój już istniejący
na jej brzegu wyznacza więc mapę $S^n\to F$.
Można ją przedłużyć na dysk dokładnie wtedy, gdy
jej klasa w $\pi_n(F)$ znika.
Dla różnych wyborów trywializacji identyfikacje
$\pi_n(F_b)$ różnią się transportem w systemie lokalnym;
zatem wszystkie te klasy są współczynnikami tego samego
kołańcucha komórkowego $o_{n+1}(s)$.

Pokażmy dokładniej, skąd pochodzi warunek kocyklu.
Trzeba uwzględnić nie tylko liczby incydencji,
lecz również działanie monodromii na $\pi_n(F)$.
Niech $e^{n+2}$ będzie komórką z mapą
charakterystyczną $\chi:D^{n+2}\to B$.
Cofnięta wiązka nad dyskiem jest trywialna.
Na podziale komórkowym $\partial D^{n+2}$ przekrój
wyznaczony na $n$-szkielecie daje elementy
$\pi_n(F)$ na granicach $(n+1)$-komórek.
Każdą z nich transportujemy do wspólnego włókna
nad wybranym wierzchołkiem $\partial D^{n+2}$.
Suma ze znakami incydencji jest dokładnie wartością
$\delta o_{n+1}(s)$ na $e^{n+2}$.

Aby wykazać jej znikanie, używamy relatywnego
kompleksu homotopii pary
$(D^{n+2},\partial D^{n+2})$ oraz naturalności
homomorfizmów brzegowych dla map charakterystycznych.
Element odpowiadający obstrukcji jest obrazem
brzegu odpowiedniej klasy względnej, a następne
zastosowanie homomorfizmu brzegowego daje zero
($\partial^2=0$ w łańcuchowym opisie komórkowym
tej pary). W lokalnej trywializacji wszystkie
transporty są identyczne, a po powrocie do $B$
odzyskujemy je przez powyższą identyfikację systemu
lokalnego. Stąd $\delta o_{n+1}(s)=0$.
Jest to \emph{komórkowy lemat o kocyklu przeszkody};
samo przytoczenie $\partial^2=0$ bez opisania
transportów byłoby niewystarczające.

Jeżeli zmienimy przekrój na komórkach wymiaru $n$,
zachowując go na $(n-1)$-szkielecie, różnice
opisują klasy map sferowych w $\pi_n(F)$; współczynniki
tej zmiany tworzą $n$-kołańcuch $b$, a nowa przeszkoda jest
$o_{n+1}(s)+\delta b$. Dlatego klasa kohomologii
nie zależy od takiej zmiany. Gdy $[o_{n+1}(s)]=0$,
wybieramy $b$ z $\delta b=-o_{n+1}(s)$,
modyfikujemy przekrój na $n$-komórkach i otrzymujemy
zerowe przeszkody dla wszystkich $(n+1)$-komórek.
Podnoszenie homotopii i rozszerzanie komórka po komórce
dają przekrój na $B^{n+1}$. Na dalszych szkieletach
problem się powtarza dla kolejnych $\pi_j(F)$.
\end{proof}

\begin{uwaga}[Zakres hipotez]
Gdy $n=1$ grupa $\pi_1(F)$ może nie być abelowa i
pierwsza przeszkoda wymaga sformułowania w teorii
nieabelowych kocykli. Dla $n\ge2$ odpowiednie grupy
są abelowe. Dla fibracji niebędących wiązkami lokalnie
trywialnymi potrzebna jest wersja komórkowej teorii
przeszkód z podnoszeniem homotopii; nie wyprowadzamy
jej tu przez pozorną lokalną trywializację.
\end{uwaga}

\begin{wniosek}[Klasa Eulera i przekrój nigdzie niezerowy]
Niech $\xi\to B$ będzie orientowaną rzeczywistą wiązką
wektorową rzędu $r\ge2$ nad CW-kompleksem $B$.
Pierwsza przeszkoda skonstruowania przekroju jednostkowej
wiązki sfer $S(\xi)\to B$ leży w $H^r(B;\Z)$ i
jest równa klasie Eulera $e(\xi)$, z ustaloną konwencją
orientacji. W szczególności $e(\xi)\ne0$ wyklucza
przekrój nigdzie niezerowy.
\end{wniosek}
\begin{proof}
Włókno $S^{r-1}$ jest $(r-2)$-spójne i
$\pi_{r-1}(S^{r-1})\cong\Z$. Dla $r\ge3$
możemy zastosować Twierdzenie~\ref{thm:obstruction-section-v6}
z $n=r-1\ge2$. Przypadek $r=2$ rozpatrujemy
osobno: zorientowana wiązka płaszczyzn redukuje
grupę strukturalną do $SO(2)\cong U(1)$,
a jej przekrój jednostkowy jest równoważny
trywializacji tej wiązki. Jej pierwszą przeszkodą
jest klasa $c_1=e\in H^2(B;\Z)$, zgodnie z klasyfikacją
wiązki liniowej zespolonej przez $BU(1)$.
Dla $r\ge3$ przeszkoda na $r$-komórce jest
stopniem mapy jej brzegu do sfery włókna. W konstrukcji Gysina ten sam stopień
otrzymujemy, ograniczając klasę Thoma do dysku
charakterystycznego i przechodząc do zwykłej
kohomologii; naturalność homomorfizmu łączącego
identyfikuje więc kołańcuch przeszkody z reprezentantem
klasy Eulera. Brak przekroju jest wówczas konieczny.
Nie twierdzimy odwrotności bez dodatkowych hipotez:
po zniknięciu klasy Eulera mogą istnieć wyższe przeszkody.
\end{proof}

\begin{przyklad}[Włosy na sferze]
Dla sfery $S^{2m}$ mamy
$\langle e(TS^{2m}),[S^{2m}]\rangle
=\chi(S^{2m})=2$ na mocy twierdzenia
Poincarégo--Hopfa. Ponieważ klasa Eulera jest niezerowa,
żadna parzystowymiarowa sfera nie ma globalnego
nigdzie niezerowego pola stycznego. Dla $S^{2m+1}$
charakterystyka Eulera znika; w tym przypadku pole
nigdzie niezerowe można dodatkowo skonstruować
jawnie, stosując na $\R^{2m+2}\cong\C^{m+1}$
wektor $x\mapsto ix$, styczny do sfery i niezerowy.
\end{przyklad}

\begin{przyklad}[Znikanie klasy Eulera nie wystarcza]
Istnieje orientowana rzeczywista wiązka wektorowa rzędu
$3$ nad $S^4$, która nie ma przekroju nigdzie niezerowego,
chociaż jej klasa Eulera jest równa zeru.
Istotnie, grupa $\pi_3(SO(3))\cong\Z$ jest niezerowa:
przykładowo dwukrotne nakrycie $S^3\to SO(3)$
przez jednostkowe kwaterniony indukuje izomorfizm
w trzeciej grupie homotopii. Wybierzmy wiązkę $\xi$
sklejoną na równiku $S^3\subset S^4$ funkcją
reprezentującą niezerowy element tej grupy.
Z klasyfikacji przez funkcje sklejające wynika,
że $\xi$ jest nietrywialna. Tymczasem
$e(\xi)\in H^3(S^4;\Z)=0$.
Gdyby $\xi$ miała przekrój nigdzie niezerowy,
rozkładałaby się jako
$\xi\cong\underline\R\oplus\eta$,
gdzie $\eta$ jest orientowaną rzeczywistą wiązką
rzędu $2$. Ale $\eta$ musi być trywialna, gdyż
$\pi_3(SO(2))=\pi_3(S^1)=0$, a więc wszystkie
takie wiązki nad $S^4$ są trywialne.
Otrzymalibyśmy sprzeczność. Jest to konkretny przykład,
w którym \emph{pierwsza} przeszkoda znika,
lecz przeszkoda wyższego rzędu pozostaje.
\end{przyklad}

\section{Funkcje sklejające i wiązki nad sferami}
\begin{twierdzenie}[Konstrukcja przez równik]
Niech $G$ będzie spójną łukowo grupą topologiczną,
a $n\ge2$. Główne wiązki $G$ nad $S^n$ powstają
przez sklejenie dwóch trywialnych wiązek nad
półkulami za pomocą funkcji
$g:S^{n-1}\to G$ na równiku. Klasy takich wiązek
odpowiadają elementom $\pi_{n-1}(G)$, z naturalną
konwencją dotyczącą punktów bazowych; dla niepunktowanych
funkcji sklejających odpowiednią relację homotopii
należy uwzględnić w klasyfikacji.
\end{twierdzenie}
\begin{proof}
Półkule $D^n_\pm$ są ściągalne, więc każda główna wiązka
nad nimi jest trywialna. Dwie trywializacje różnią się
na przekroju równikowym przez mapę do $G$.
Sklejając $(D^n_+\times G)$ z $(D^n_-\times G)$
według relacji $(x,h)_+\sim(x,g(x)h)_-$ otrzymujemy
wiązkę. Zmiana trywializacji nad półkulą działa na
g przez mnożenie z lewej lub prawej strony funkcją,
która rozszerza się na dysk, a zatem jest homotopijna
do stałej. Gdy $G$ jest spójna łukowo, odpowiednie
stałe są połączone drogami i otrzymujemy klasyfikację
przez homotopię funkcji z równika do $G$.
Dla $n\ge2$ i spójnej łukowo grupy $G$ jest to
$\pi_{n-1}(G)$ (w szczególności dla $n=2$ jest to
abelowa grupa $\pi_1(G)$).
\end{proof}
\begin{przyklad}
Dla $G=U(1)=S^1$ oraz $n=2$ funkcja
$g_k(e^{it})=e^{ikt}$ ma stopień $k$.
Sklejona zespolona wiązka liniowa ma
$\langle c_1(L_k),[S^2]\rangle=k$ z przyjętą
orientacją i konwencją klejenia. Dla $k=0$ jest trywialna,
a dla $k=\pm1$ otrzymujemy tautologiczną wiązkę
nad $\C P^1$ albo jej dualną (zależnie od znaku).
\end{przyklad}


\section{Uzupełnienia audytu: konstrukcje i granice hipotez}
\label{sec:audyt-VI-VII}
\begin{lemat}[Komórki przy konstrukcji $K(G,n)$]
Niech $A$ będzie CW-kompleksem i niech $Y=A\cup_f e^{m}$
dla $m\ge3$. Dla każdego $j\le m-2$ włączenie
$A\hookrightarrow Y$ indukuje izomorfizm
$\pi_j(A)\to\pi_j(Y)$. Dla $j=m-1$ indukuje
surjekcję, której jądro jest normalnie generowane
przez klasę mapy dołączającej (w wymiarach $m-1\ge2$
z naturalnym uwzględnieniem działania $\pi_1$).
\end{lemat}
\begin{proof}
Para $(Y,A)$ jest, przez przybliżenie komórkowe
względem $A$, $(m-1)$-spójna: klasy względne w stopniach
$m-1$ i niższych można reprezentować mapami dysków,
których obraz po przybliżeniu komórkowym unika wnętrza
jedynej nowej $m$-komórki. W długim ciągu
$\pi_j(A)\to\pi_j(Y)\to\pi_j(Y,A)$
otrzymujemy zatem izomorfizmy do wymiaru $m-2$.
W stopniu $m-1$ jądro mapy włączenia jest obrazem
$\partial:\pi_m(Y,A)\to\pi_{m-1}(A)$,
a relatywny generator ma jako brzeg klasę mapy $f$.
Dla kilku komórek stosujemy argument kolejno,
a następnie przechodzimy do CW-granicy; zwartość
sfer gwarantuje faktoryzację map przez pewien etap.
\end{proof}

\begin{uwaga}[Co jest rzeczywiście udowodnione]
Lemat wyjaśnia konstrukcję istnienia $K(G,n)$.
W klasyfikacji \emph{wszystkich} wiązek głównych
występuje oddzielne twierdzenie o uniwersalnym
numerowalnym modelu $EG\to BG$: jego użycie wymaga
parazwartości i własności Hausdorffa podstawy
oraz kontrolowania lokalnie skończonych podziałów
jedności. Wzór na mapę klasyfikującą pokazuje
surjektywność klasyfikacji i jej naturalność;
pełna homotopijna jednoznaczność zawiera dodatkowo
argument o wiązkach nad cylindrem $X\times I$.
Nie zastępujemy tego argumentu samym stwierdzeniem,
że dwie mapy są punktowo podobne.
\end{uwaga}

\begin{przyklad}[Rozszerzenie na komórkę w dwóch etapach]
Rozważmy przestrzeń Moore'a
$M(\Z/m,n)=S^n\cup_{m}e^{n+1}$,
$m\ge2$, $n\ge2$. Mapa $S^n\to K(A,n)$
reprezentująca $a\in A$ przedłuża się na
$M(\Z/m,n)$ dokładnie wtedy, gdy $ma=0$.
Istotnie, po złożeniu z mapą dołączającą stopnia $m$
otrzymujemy $ma\in\pi_n(K(A,n))=A$, a przedłużenie
istnieje dokładnie dla klasy zerowej. Wynik jest zgodny
z $H^n(M(\Z/m,n);A)\cong A[m]$.
\end{przyklad}

\section{Podsumowanie i granice obecnej iteracji}
Uzyskaliśmy trzy mechanizmy służące różnym celom:
$K(A,n)$ reprezentuje kohomologię,
$BG$ klasyfikuje wiązki, a kołańcuchy przeszkód
wykrywają możliwość rozszerzania map i przekrojów.
Użyte konstrukcje są zgodne z poprzednim rozdziałem:
przestrzenie klasyfikujące i klasy Eulera można badać
ciągami spektralnymi fibracji, a pierwsza przeszkoda
jest obliczana w kohomologii z lokalnymi współczynnikami.

Wieże Postnikowa, klasy $k$-niezmienników,
operacje Steenroda i szczegółowa teoria przeszkód
z wieloma niezerowymi grupami homotopii włókna
stanowią naturalny następny etap. Nie utożsamiamy
pierwszej przeszkody z całą teorią przeszkód.


% Rozdzial glowny iteracji VII. Do wstawienia przed \appendix.
